\section{宏观经济数据}\label{ux5b8fux89c2ux7ecfux6d4eux6570ux636e}

\subsection{GDP}\label{gdp}

Two definitions:

\begin{itemize}
\tightlist
\item
  Total \textbf{expenditure} on domestically produced final goods and services.
\item
  Total \textbf{income} earned by domestically located factors of production.
\end{itemize}

\begin{quote}
Expenditure = Income.
Since every dollar a buyer spends becomes income to the seller.
\end{quote}

\subsection{Measures}\label{measures}

\subsubsection{Value added}\label{value-added}

\textbf{Value added:} The value of output minus the value of the intermediate goods used to produce that output.

\[
\text{silica }(1\,\mathrm{RMB}) \longrightarrow \text{glass }(4\,\mathrm{RMB}) \longrightarrow \text{automobile }(100\,\mathrm{RMB})
\]

VA: \(1 - 0 = 1, 4 - 1 = 3, 100 - 4 = 96\).\\
GDP: \(\sum \text{VA} = 1 + 3 + 96 = 100\).

\begin{quote}
GDP = value of final goods produced = sum of value added at all stages of production.
\end{quote}

\subsubsection{The expenditure components of GDP}\label{the-expenditure-components-of-gdp}

\[
Y = C + I + G + NX
\]
\(C\): consumption, \(I\): investment, \(G\): government spending, \(NX\): net exports.

\(Y\): value of total output, \(C, I, G, NX\): aggregate expenditure.

\paragraph{Consumption}\label{consumption}

The value of all goods and services bought by households.\\
\textbf{Durable goods:} Last a long time.\\
\textbf{Nondurable goods:} Last a short time.\\
\textbf{Services:} Intangible/non-physical items or activities.

\paragraph{Investment}\label{investment}

Spending on capital, a physical asset used in future production.\\
\textbf{Business fixed investment:} Plant / Equipment\\
\textbf{Residential fixed investment:} Housing units.\\
\textbf{Inventory investment:} \textbf{\emph{The change}} in the value of all firms' inventories.

\begin{quote}
Investment is spending on new capital.\\
即是说，一个是流量，一个是存量。\\
\textbf{Flow}: a quantity measured per unit of time.\\
\textbf{Stock}: a quantity measured at a point in time.
\end{quote}

\paragraph{Government Spending}\label{government-spending}

Including all government spending on goods and services.\\
\textbf{Excluding} transfer payments. 没有转移支付.

\paragraph{Net exports}\label{net-exports}

\textbf{NX} = exports - imports.

\begin{itemize}
\tightlist
\item
  Exports: The value of goods and services sold to other countries.
\item
  Imports: The value of goods and services purchased from other countries.
\end{itemize}

\subsection{Components of National Income（国民收入构成）}\label{components-of-national-incomeux56fdux6c11ux6536ux5165ux6784ux6210}

讲义比较了 2014 年与 2022 年美国 \textbf{National Income} 的构成：

\begin{longtable}[]{@{}lrr@{}}
\toprule\noalign{}
Component & 2014 & 2022 \\
\midrule\noalign{}
\endhead
\bottomrule\noalign{}
\endlastfoot
Compensation of employees（雇员报酬） & 61\% & 62\% \\
Proprietors' income（业主收入） & 9\% & 8\% \\
Rental income（租金收入） & 4\% & 4\% \\
Corporate profits（公司利润） & 14\% & 14\% \\
Net interest（净利息） & 4\% & 2\% \\
Indirect business taxes and other（间接营业税及其他） & 8\% & 10\% \\
\end{longtable}

\textbf{Key observation}

雇员报酬占比约六成，是国民收入中最大的部分。两年间总体结构较稳定，较明显的变化是净利息占比下降、间接营业税及其他项目占比上升。

\subsection{Other Measures of Income（其他收入指标）}\label{other-measures-of-incomeux5176ux4ed6ux6536ux5165ux6307ux6807}

从 \textbf{National Income} 到家庭实际可支配资源，需要继续扣除``赚得但未直接收到''的项目，并加回``收到但不属于当期要素收入''的项目：

\[
\begin{aligned}
\text{Personal Income}
={}&\ \text{National Income}
-\text{Indirect Business Taxes}\\
&-\text{Corporate Profits}
-\text{Social Insurance Contributions}
-\text{Net Interest}\\
&+\text{Dividends}+\text{Government Transfers to Individuals}\\
&+\text{Personal Interest Income}.
\end{aligned}
\]

个人可支配收入为

\[
\text{Disposable Personal Income}
=\text{Personal Income}
-\text{Personal Tax and Nontax Payments}.
\]

\textbf{Interpretation}

\textbf{Disposable Personal Income, DPI} 是家庭和非公司制企业最终能够用于消费或储蓄的收入：

\[
\text{DPI}=\text{Consumption}+\text{Personal Saving}.
\]

\subsection{Nominal GDP and Real GDP（名义与实际 GDP）}\label{nominal-gdp-and-real-gdpux540dux4e49ux4e0eux5b9eux9645-gdp}

\textbf{Nominal GDP} 使用当期价格评价当期产量，因此价格与数量变化都会影响其数值。\textbf{Real GDP} 使用某个基年的价格评价各期产量，从而控制价格变化，只反映产出数量的变化。

设经济中有 \(n\) 种最终产品，\(P_{it}\) 和 \(Q_{it}\) 分别表示产品 \(i\) 在 \(t\) 期的价格和数量，\(B\) 表示基年，则

\[
GDP_t=\sum_{i=1}^{n}P_{it}Q_{it},
\]

\[
RGDP_t^{(B)}=\sum_{i=1}^{n}P_{iB}Q_{it}.
\]

\begin{longtable}[]{@{}llll@{}}
\toprule\noalign{}
Measure & Price used & Quantity used & What can change it? \\
\midrule\noalign{}
\endhead
\bottomrule\noalign{}
\endlastfoot
Nominal GDP & 当期价格 \(P_{it}\) & 当期数量 \(Q_{it}\) & 价格或数量 \\
Real GDP & 固定基年价格 \(P_{iB}\) & 当期数量 \(Q_{it}\) & 仅数量 \\
\end{longtable}

\begin{quote}
比较不同时期的生产活动时，应使用 real GDP；nominal GDP 的增长不能直接解释为实际产出增长。
\end{quote}

\subsection{Fixed Base-Year Growth（固定基年增长率）}\label{fixed-base-year-growthux56faux5b9aux57faux5e74ux589eux957fux7387}

固定基年价格下，\(t-1\) 到 \(t\) 期的实际 GDP 增长因子为

\[
1+g_t^{(B)}
=\frac{RGDP_t^{(B)}}{RGDP_{t-1}^{(B)}}
=\frac{\sum_{i=1}^{n}P_{iB}Q_{it}}
{\sum_{i=1}^{n}P_{iB}Q_{i,t-1}}.
\]

把分子中的每一项写成

\[
P_{iB}Q_{it}
=P_{iB}Q_{i,t-1}\left(\frac{Q_{it}}{Q_{i,t-1}}\right),
\]

可得

\[
1+g_t^{(B)}
=\sum_{i=1}^{n}\omega_{iB,t-1}
\left(\frac{Q_{it}}{Q_{i,t-1}}\right),
\]

其中

\[
\omega_{iB,t-1}
=\frac{P_{iB}Q_{i,t-1}}
{\sum_{j=1}^{n}P_{jB}Q_{j,t-1}},
\qquad
\sum_{i=1}^{n}\omega_{iB,t-1}=1.
\]

\textbf{Interpretation}

整体实际增长因子是各产品数量增长因子的加权平均；权重是各产品按基年价格计算的上期产值占比。

\subsection{Why a Fixed Base Year Is Problematic（固定基年的问题）}\label{why-a-fixed-base-year-is-problematicux56faux5b9aux57faux5e74ux7684ux95eeux9898}

相对价格会随时间改变，因此固定基年法得到的增长率取决于基年 \(B\) 的选择。

\begin{itemize}
\tightlist
\item
  若某行业数量增长快、相对价格持续下降，较早基年的高价格会使该行业被赋予过大权重；
\item
  若某行业数量增长慢，则可能出现相反方向的权重偏差；
\item
  基年距离当前越远，固定权重越难代表当前经济结构。
\end{itemize}

因此，同一组数量变化可能因基年不同而产生不同的 measured growth。

\subsection{Chain-Weighted Real GDP（链式加权实际 GDP）}\label{chain-weighted-real-gdpux94feux5f0fux52a0ux6743ux5b9eux9645-gdp}

\textbf{Chain weighting} 每年更新价格基准，并结合相邻两年的价格权重，减少固定基年过时造成的偏差。讲义指出，这是美国 \textbf{Bureau of Economic Analysis, BEA} 当前采用的官方实际 GDP 度量方法。

\subsubsection{Step 1：使用上一期价格}\label{step-1ux4f7fux7528ux4e0aux4e00ux671fux4ef7ux683c}

用 \(t-1\) 期价格评价从 \(t-1\) 到 \(t\) 的数量变化：

\[
(1+g_t)_{t-1}
=\frac{\sum_{i=1}^{n}P_{i,t-1}Q_{it}}
{\sum_{i=1}^{n}P_{i,t-1}Q_{i,t-1}}
=\sum_{i=1}^{n}\omega_{i,t-1}
\left(\frac{Q_{it}}{Q_{i,t-1}}\right),
\]

其中

\[
\omega_{i,t-1}
=\frac{P_{i,t-1}Q_{i,t-1}}
{\sum_{j=1}^{n}P_{j,t-1}Q_{j,t-1}}.
\]

\subsubsection{Step 2：使用当期价格}\label{step-2ux4f7fux7528ux5f53ux671fux4ef7ux683c}

再用 \(t\) 期价格评价同一数量变化：

\[
(1+g_t)_t
=\frac{\sum_{i=1}^{n}P_{it}Q_{it}}
{\sum_{i=1}^{n}P_{it}Q_{i,t-1}}
=\sum_{i=1}^{n}\omega_{i,t}
\left(\frac{Q_{it}}{Q_{i,t-1}}\right),
\]

其中

\[
\omega_{i,t}
=\frac{P_{it}Q_{i,t-1}}
{\sum_{j=1}^{n}P_{jt}Q_{j,t-1}}.
\]

截至 Page 42，我们得到了同一数量增长的两种估计：\((1+g_t)_{t-1}\) 与 \((1+g_t)_t\)。后续页面将把二者合成为 chain-weighted growth rate。

\begin{enumerate}
\def\labelenumi{\arabic{enumi}.}
\tightlist
\item
  \textbf{Personal Income} 从 national income 调整到个人实际收到的收入；\textbf{DPI} 再扣除个人税费。
\item
  \textbf{Nominal GDP} 混合了价格与数量变化；\textbf{Real GDP} 隔离数量变化。
\item
  固定基年增长率本质上是产品数量增长因子的加权平均。
\item
  固定权重会随时间失真，因此 chain weighting 使用相邻年份的价格信息持续更新权重。
\end{enumerate}

上一节分别使用 \(t-1\) 期价格和 \(t\) 期价格计算了同一数量变化：

\[
(1+g_t)_{t-1},
\qquad
(1+g_t)_t.
\]

\subsubsection{Step 3：合成相邻两年的权重}\label{step-3ux5408ux6210ux76f8ux90bbux4e24ux5e74ux7684ux6743ux91cd}

Chain-weighted growth factor 取两种增长因子的几何平均：

\[
1+g_t
=\left[(1+g_t)_t(1+g_t)_{t-1}\right]^{1/2}.
\]

几何平均对两个方向的比较保持对称：既不会只依赖上一期价格，也不会只依赖当期价格。

\subsubsection{From growth rates to levels（从增长率还原水平值）}\label{from-growth-rates-to-levelsux4eceux589eux957fux7387ux8fd8ux539fux6c34ux5e73ux503c}

先选择参考年 \(r\)，令该年实际 GDP 等于名义 GDP，再把链式增长率逐年连接：

\[
RGDP_r^{(r)}=GDP_r,
\]

\[
RGDP_t^{(r)}
=GDP_r\prod_{s=r+1}^{t}(1+g_s).
\]

讲义以 \(r=t-4\) 为例：

\[
RGDP_t
=(1+g_t)(1+g_{t-1})(1+g_{t-2})(1+g_{t-3})GDP_{t-4}.
\]

所得实际 GDP 以参考年 \(r\) 的货币价值表示。参考年只决定水平值的单位；各期 chain-weighted growth rates 不因固定基年的选择而改变。

\subsection{GDP Deflator（GDP 平减指数）}\label{gdp-deflatorgdp-ux5e73ux51cfux6307ux6570}

\textbf{Inflation rate} 是总体价格水平的百分比上升。GDP deflator 是衡量总体价格水平的一种指标：

\[
\text{GDP Deflator}_t
=100\times\frac{\text{Nominal GDP}_t}{\text{Real GDP}_t}.
\]

相应的通胀率为

\[
\pi_t
=\frac{\text{Deflator}_t-\text{Deflator}_{t-1}}
{\text{Deflator}_{t-1}}\times100\%.
\]

\subsubsection{Fixed-weight representation（固定权重表示）}\label{fixed-weight-representationux56faux5b9aux6743ux91cdux8868ux793a}

若实际 GDP 使用基年 \(B\) 的固定价格，记

\[
NGDP_t=\sum_{i=1}^{n}P_{it}Q_{it},
\qquad
RGDP_t^{(B)}=\sum_{i=1}^{n}P_{iB}Q_{it},
\]

则

\[
\frac{\text{Deflator}_t}{100}
=\frac{\sum_{i=1}^{n}P_{it}Q_{it}}
{\sum_{i=1}^{n}P_{iB}Q_{it}}.
\]

把 \(P_{it}\) 写成 \(P_{iB}(P_{it}/P_{iB})\)，可得

\[
\frac{\text{Deflator}_t}{100}
=\sum_{i=1}^{n}\gamma_{it}
\left(\frac{P_{it}}{P_{iB}}\right),
\]

其中

\[
\gamma_{it}
=\frac{P_{iB}Q_{it}}
{\sum_{j=1}^{n}P_{jB}Q_{jt}},
\qquad
\sum_{i=1}^{n}\gamma_{it}=1.
\]

因此 GDP deflator 是各产品 \textbf{price relative} \(P_{it}/P_{iB}\) 的加权平均。权重包含当期数量 \(Q_{it}\)，会随当前生产结构变化，所以在固定基年 real GDP 下，它表现为 \textbf{Paasche index}（当期数量加权价格指数）。

\subsubsection{Chain-weighted deflator（链式加权平减指数）}\label{chain-weighted-deflatorux94feux5f0fux52a0ux6743ux5e73ux51cfux6307ux6570}

若 real GDP 采用 chain weighting，也可以用同样的``nominal / real''方式构造 deflator。其数量权重逐年更新，并综合当年与上年的数量信息，避免价格指数长期依赖过时的固定篮子。

\subsection{Percentage-Change Approximations（百分比变化近似）}\label{percentage-change-approximationsux767eux5206ux6bd4ux53d8ux5316ux8fd1ux4f3c}

当百分比变化不大时，有两个常用近似。

\subsubsection{Product rule（乘积法则）}\label{product-ruleux4e58ux79efux6cd5ux5219}

\[
\%\Delta(XY)
\approx\%\Delta X+\%\Delta Y.
\]

例如，时薪上升 \(5\%\)、工作时长上升 \(7\%\)，则工资收入近似上升

\[
5\%+7\%=12\%.
\]

精确增长率为 \((1.05)(1.07)-1=12.35\%\)；近似式忽略了两个变化率的乘积项。

\subsubsection{Quotient rule（商法则）}\label{quotient-ruleux5546ux6cd5ux5219}

\[
\%\Delta\left(\frac XY\right)
\approx\%\Delta X-\%\Delta Y.
\]

由于

\[
\text{GDP Deflator}=100\times\frac{NGDP}{RGDP},
\]

若 nominal GDP 上升 \(9\%\)、real GDP 上升 \(4\%\)，则通胀率近似为

\[
9\%-4\%=5\%.
\]

\subsection{Identifying Recessions（如何判定经济衰退）}\label{identifying-recessionsux5982ux4f55ux5224ux5b9aux7ecfux6d4eux8870ux9000}

常见经验法则是``real GDP 连续两个季度下降''，但这不是完整的官方判定标准。

\textbf{NBER approach}

NBER 的商业周期判定更强调：经济活动显著下降、影响范围广、持续数月以上，并综合观察月度指标，例如：

\begin{itemize}
\tightlist
\item
  industrial production；
\item
  employment；
\item
  real income；
\item
  wholesale and retail trade。
\end{itemize}

因此，recession 是多维经济活动的共同收缩，而不只是单个季度 GDP 数字触发的机械标签。

\subsection{Consumer Price Index（消费者价格指数）}\label{consumer-price-indexux6d88ux8d39ux8005ux4ef7ux683cux6307ux6570}

\textbf{CPI} 衡量典型家庭购买的一篮子商品与服务的总体价格水平，由 \textbf{Bureau of Labor Statistics, BLS} 发布。

主要用途包括：

\begin{itemize}
\tightlist
\item
  跟踪典型家庭生活成本的变化；
\item
  对工资、养老金等合同进行 cost-of-living adjustment（COLA）；
\item
  把不同时期的货币金额转换到可比价格水平。
\end{itemize}

\subsubsection{Construction（构造过程）}\label{constructionux6784ux9020ux8fc7ux7a0b}

\begin{enumerate}
\def\labelenumi{\arabic{enumi}.}
\tightlist
\item
  调查消费者，确定典型消费篮子的组成；
\item
  每月收集篮子中各项目的价格并计算篮子总成本；
\item
  用当月成本与基期成本之比构造 CPI。
\end{enumerate}

设 \(C_i\) 为篮子中商品 \(i\) 的固定数量，\(P_{it}\) 为其在 \(t\) 期的价格，则

\[
E_t=\sum_{i=1}^{n}C_iP_{it},
\qquad
E_B=\sum_{i=1}^{n}C_iP_{iB},
\]

\[
CPI_t=100\times\frac{E_t}{E_B}.
\]

进一步可写为

\[
\frac{CPI_t}{100}
=\sum_{i=1}^{n}\gamma_{iB}
\left(\frac{P_{it}}{P_{iB}}\right),
\]

其中基期权重为

\[
\gamma_{iB}
=\frac{C_iP_{iB}}
{\sum_{j=1}^{n}C_jP_{jB}},
\qquad
\sum_{i=1}^{n}\gamma_{iB}=1.
\]

CPI 是各商品价格相对数的加权平均；权重等于该商品在基期篮子支出中的占比。由于篮子数量固定，CPI 属于 \textbf{Laspeyres index}。

\subsubsection{Why CPI may overstate inflation（CPI 的向上偏差）}\label{why-cpi-may-overstate-inflationcpi-ux7684ux5411ux4e0aux504fux5dee}

\begin{enumerate}
\def\labelenumi{\arabic{enumi}.}
\tightlist
\item
  \textbf{Substitution bias}：固定篮子无法及时反映消费者转向相对价格下降商品的行为；
\item
  \textbf{New goods}：新产品扩大选择并提高货币的实际购买价值，但固定篮子不能立刻体现这部分福利改善；
\item
  \textbf{Unmeasured quality change}：质量提高意味着同样货币购买到更高价值，但质量调整往往不完全。
\end{enumerate}

讲义指出，1995 年专家委员会估计 CPI 每年高估通胀约 \(1.1\) 个百分点；BLS 随后进行了修正，目前讲义给出的判断是偏差可能低于每年 \(1\%\)。

\subsection{Comparing Price Indexes（价格指数比较）}\label{comparing-price-indexesux4ef7ux683cux6307ux6570ux6bd4ux8f83}

\subsubsection{CPI vs.~GDP deflator}\label{cpi-vs.-gdp-deflator}

\begin{longtable}[]{@{}>{\raggedright\arraybackslash}p{0.22\linewidth}>{\raggedright\arraybackslash}p{0.34\linewidth}>{\raggedright\arraybackslash}p{0.34\linewidth}@{}}
\toprule\noalign{}
Dimension & CPI & GDP deflator \\
\midrule\noalign{}
\endhead
\bottomrule\noalign{}
\endlastfoot
Coverage & 家庭消费的商品与服务 & 国内生产的全部最终产品与服务 \\
Imported consumer goods & 包含 & 不包含，因为不属于国内生产 \\
Domestically produced capital goods & 不包含 & 包含 \\
Basket / weights & 固定消费篮子 & 随当前国内生产结构变化 \\
Index type in the fixed-base setup & Laspeyres & Paasche \\
\end{longtable}

因此，两者衡量的``总体价格''并不相同：CPI 更接近消费者生活成本，GDP deflator 更接近国内总产出的价格。

\subsubsection{PCE deflator}\label{pce-deflator}

\textbf{Personal Consumption Expenditures Price Index} 是 nominal consumer spending 与 real consumer spending 的比率。

它与 CPI 的共同点：

\begin{itemize}
\tightlist
\item
  只覆盖消费支出；
\item
  包含进口消费品。
\end{itemize}

它与 GDP deflator 的共同点：

\begin{itemize}
\tightlist
\item
  篮子与权重会随时间改变，能够反映消费替代。
\end{itemize}

讲义指出，Federal Reserve 更偏好 PCE price index。

\subsubsection{Core inflation（核心通胀）}\label{core-inflationux6838ux5fc3ux901aux80c0}

BLS 与 BEA 都会发布排除 food and energy 的 CPI/PCE 指标，称为 \textbf{core inflation}。食品与能源价格短期波动较大，排除它们有助于观察更持久的 underlying price trend，但 core inflation 并不代表家庭不需要购买这些项目。

\subsection{Labor-Market Statistics（劳动力市场统计）}\label{labor-market-statisticsux52b3ux52a8ux529bux5e02ux573aux7edfux8ba1}

\subsubsection{Population categories（人口分类）}\label{population-categoriesux4ebaux53e3ux5206ux7c7b}

\begin{longtable}[]{@{}ll@{}}
\toprule\noalign{}
Category & Definition \\
\midrule\noalign{}
\endhead
\bottomrule\noalign{}
\endlastfoot
Employed, \(E\) & 正在从事有报酬工作 \\
Unemployed, \(U\) & 没有工作，但正在寻找工作 \\
Labor force, \(L\) & 所有就业者与失业者，\(L=E+U\) \\
Not in the labor force & 没有工作且没有寻找工作 \\
\end{longtable}

``没有工作''不自动等于 unemployed；只有正在求职者才计入 \(U\)，其余没有工作者属于 not in the labor force。

\subsubsection{Two key rates（两个核心比率）}\label{two-key-ratesux4e24ux4e2aux6838ux5fc3ux6bd4ux7387}

失业率为劳动力中失业者的比例：

\[
u=\frac{U}{L}\times100\%.
\]

劳动力参与率为成年人口 \(POP\) 中进入劳动力市场的比例：

\[
LFPR=\frac{L}{POP}\times100\%.
\]

BLS 主要通过 \textbf{Current Population Survey, CPS}（家庭调查）测量这些概念。

\subsubsection{Example：2021 年 2 月美国劳动力数据}\label{example2021-ux5e74-2-ux6708ux7f8eux56fdux52b3ux52a8ux529bux6570ux636e}

给定

\[
E=150.24,\qquad U=9.72,\qquad POP=260.66
\]

（单位均为 million），则

\[
L=E+U=150.24+9.72=159.96,
\]

\[
u=\frac{9.72}{159.96}\times100\%\approx6.1\%,
\]

\[
LFPR=\frac{159.96}{260.66}\times100\%\approx61.37\%.
\]

\subsubsection{Percentage changes in labor ratios（劳动力比率的百分比变化）}\label{percentage-changes-in-labor-ratiosux52b3ux52a8ux529bux6bd4ux7387ux7684ux767eux5206ux6bd4ux53d8ux5316}

若人口增加 \(1\%\)、劳动力增加 \(3\%\)、失业人数增加 \(2\%\)，利用商的百分比变化近似：

\[
\%\Delta LFPR
\approx\%\Delta L-\%\Delta POP
=3\%-1\%=2\%,
\]

\[
\%\Delta u
\approx\%\Delta U-\%\Delta L
=2\%-3\%=-1\%.
\]

\textbf{Warning}

这里得到的是相对于原比率的百分比变化，不是 percentage-point change。例如初始 \(LFPR=60.0\%\)，增加 \(2\%\) 后为

\[
60.0\%\times1.02=61.2\%,
\]

即增加 \(1.2\) 个百分点，而不是 \(2\) 个百分点。

\subsubsection{Household vs.~establishment survey（家庭调查与机构调查）}\label{household-vs.-establishment-surveyux5bb6ux5eadux8c03ux67e5ux4e0eux673aux6784ux8c03ux67e5}

BLS 还通过 \textbf{establishment survey} 调查企业 payrolls，得到另一套就业指标。两套调查偶尔会发生分歧，原因包括：

\begin{itemize}
\tightlist
\item
  self-employed persons 在两套调查中的处理不同；
\item
  新成立企业可能尚未进入机构调查样本；
\item
  家庭调查需要从样本推断总体人口，存在相应的统计与技术问题。
\end{itemize}

家庭调查以``人''为单位并支持就业状态、失业率和参与率的计算；机构调查以``工资岗位''为单位，通常更适合观察 nonfarm payroll employment 的月度变化。

\subsection{Chapter 2 Review}\label{chapter-2-review}

\begin{enumerate}
\def\labelenumi{\arabic{enumi}.}
\tightlist
\item
  GDP 同时衡量国内最终产出的总支出与生产要素获得的总收入，并满足 \(Y=C+I+G+NX\)。
\item
  Nominal GDP 使用当期价格；real GDP 控制价格变化并衡量实际产出。
\item
  Chain weighting 使用相邻年份的价格信息构造增长率，减少固定基年过时造成的偏差。
\item
  GDP deflator、CPI 与 PCE deflator 的覆盖范围和权重不同，回答的是不同的价格问题。
\item
  衰退判断需要综合产出、就业、工业生产、收入和销售等指标，而非只看单一规则。
\item
  失业率的分母是 labor force；没有工作但不求职者不属于 unemployed。
\item
  Unemployment rate 与 LFPR 应结合阅读，才能完整描述劳动力市场。
\end{enumerate}

